\scanpage{4}
% source line 1: title
\uline{La Théorie \unsure{du} corps de classes}

% source lines 2-4
$K$ corps de \unsure{nbs|nbres} algébriques, ${\unsure{J|\mathscr{J}}}_K$ \unsure{groupe des} idèles, $U_K$ \deleted{\unsure{unités}} sous-groupe des idèles unités (à l'infini, \unsure{ils sont} $>0$ aux places réelles), ${\unsure{D|\mathscr{D}}}^a_K$ les diviseurs absolus, $D_K$ les diviseurs $(D^a_K\approx D_K\times Z_2^{r_1})$

% source line 4a: small interlinear addition
\inserted{Soit $\underline m$ un diviseur absolu, positif aux places finies (i.e. $\underline m_0$ est un idéal \uline{entier})}

% source lines 5-13
$D^{\underline m}_K$ les diviseurs premiers à $\underline m$ (ne dépend que de la partie finie $\underline m_0$ de l'idéal absolu $\underline m$), $P^{\underline m}_K\unsure{\subset|\in}D^{\underline m}_K$ \unsure{formé} des $(x)$ où $x\in K^*$ est $\equiv1\ (\underline m)$, i.e. $v_{\mathfrak p}(x\deleted{\unsure{+}}-1)\geq v_{\mathfrak p}(\underline m)$ aux places finies où $v_{\mathfrak p}(\underline m)>0$, et \deleted{$\unsure{x}$} $x\in K_{\mathfrak p}^{*+}$ aux places infinies où $v_{\mathfrak p}(\underline m)\neq0$. Les $x$ en question forment $K^{*(\underline m)}$. ${\unsure{\mathscr{P}|\mathfrak{P}|\mathcal{P}}}^{\underline m}_K=D^{\underline m}_K/P^{\underline m}_K$. (N.B. $K_{\deleted{\unsure{\mathfrak p}}}^{*\underline m}$ est aussi formé des quotients $x/y$, où $x$ et $y$ sont \uline{entiers} tels que $v_{\mathfrak p}(x-1)\geq v_{\mathfrak p}(\underline m)$ si $\mathfrak p$ est finie, et $x\in K_{\mathfrak p}^{*+}$ si $\mathfrak p$ est infinie et $v_{\mathfrak p}(\underline m)\neq0$. On a une suite exacte
\diagram{\[
\begin{tikzcd}[ampersand replacement=\&]
\&\&\&\&\text{classes d'idéaux}\\
0\arrow[r]\&N_K^{\underline m}\arrow[r]\&\mathscr{P}_K^{\underline m}\arrow[r]\&\mathscr{P}_K\arrow[r]\arrow[ur]\&0
\end{tikzcd}
\tag*{(1)}
\]}
où le noyau $N_K^{\underline m}$ se décrit ainsi
\[
0\longrightarrow Z_2^{S_\infty}/\operatorname{Im}E_K^{\underline m_0}\longrightarrow N_K^{\underline m}\longrightarrow(A_K/\underline m_0^{\deleted{\unsure{*}}})^*/\operatorname{Im}E_K\longrightarrow0
\tag*{(2)}
\]

% source lines 17-23
($A_K$, entiers de $K$, $E_K$ unités, $E_K^{\underline m_0}$ unités qui sont $\unsure{\approx|\equiv}1\ (\underline m_0)$), $S$ \unsure{ens} des places qui figurent dans $\underline m$, $S_\infty$ la partie infinie; N.B. Si $\underline m_0$ est grand, alors $E_K^{\underline m_0}$ est contenu dans n'importe quel sous-groupe d'indice fini \inserted{de $E_K$}, donc $\operatorname{Im}E_K^{\underline m_0}$ dans $Z_2^\infty$ est nul, donc $Z_2^\infty\unsure{\in|\subset}N_K^{\underline m}$ !). Si $\underline m$ est fini, i.e. $S_\infty=\varnothing$, alors on a la suite exacte
\[
0\longrightarrow(A_K/\underline m_0)^*/\operatorname{Im}E_K\longrightarrow\mathscr{P}_K^{\underline m}\longrightarrow\mathscr{P}_K\longrightarrow0.
\]

% source lines 24-25
Soit $C_K=J_K/K^*$, $U_K^{\deleted{\unsure{S}}\underline m}=\displaystyle\prod_{\mathfrak p\in\complement S}K_{\mathfrak p}^*\times\prod_{\mathfrak p\in S_0}(\deleted{\unsure{K_{\mathfrak p}}}1+\mathfrak p^{v_{\mathfrak p}(\underline m)})\times\prod_{\mathfrak p\in S_\infty}K_{\mathfrak p}^{*+}$, i.e. \unsure{c'est} l'\unsure{ens} des \uline{idèles} qui sont $\equiv1\ (\underline m)$.

% source lines 26-29
Alors $U_K^{\unsure{\complement S|\underline m}}$ est un voisinage \unsure{ouvert} de $0$ dans $J_K$ (et si $\underline m$ varie, on obtient un \unsure{syst.} fondamental de tels \unsure{sous-gps}), et \deleted{$\unsure{C_K/\operatorname{Im}U_K^{\underline m}}$} on a
\[
C_K/\operatorname{Im}U_K^{\underline m}\approx\mathscr{P}_K^{\underline m}.
\]
% archive sheet number: 3
